\documentclass[12pt]{article}
\usepackage[T1]{fontenc}
\usepackage[utf8]{inputenc}
\usepackage{lmodern}
\usepackage{textcomp}
\usepackage{enumitem}
\usepackage{geometry}
\geometry{margin=1in}
\begin{document}
\begin{center}
Answer Key - Economics - Graduate Level Assessment 9 \
\small (Answers are ordered exactly as in the question paper.)
\end{center}

\section*{Multiple Choice}
\begin{enumerate}[label=\arabic*.]
\item \textbf{Answer:} A
\\ \textit{Options:} A) tariffs protecting domestic jobs are eliminated.; B) the literacy rate improves.; C) the tax rate is lowered.; D) the inflation rate decreases.; E) there is an increase in technological advancements.
\\ \textit{Note:} Eliminating tariffs protecting domestic jobs can lead to increased competition and efficiency in the economy, which may result in a reallocation of resources that ultimately reduces overall potential GDP ceteris paribus, as some industries may contract while others grow.
\item \textbf{Answer:} A
\\ \textit{Options:} A) Conducting monetary policy for the government; B) Offering educational services on finance; C) Investing in securities; D) Issuing credit cards; E) Setting interest rates for the entire economy
\\ \textit{Note:} The correct answer is justified because commercial banks play a critical role in implementing the government's monetary policy by influencing money supply, interest rates, and overall economic stability through their lending and deposit activities.
\item \textbf{Answer:} B
\\ \textit{Options:} A) an increase in net exports; B) a decrease in net exports; C) an increase in the real interest rate; D) an increase in the nominal interest rate
\\ \textit{Note:} An easy money (expansionary) policy lowers interest rates, which can lead to a depreciation of the domestic currency, making exports cheaper and imports more expensive; however, if the currency appreciates instead due to increased capital inflows, net exports can decrease as domestic goods become relatively more expensive for foreign buyers.
\item \textbf{Answer:} A
\\ \textit{Options:} A) Interest is the penalty for early consumption; B) Interest is the reward for immediate consumption; C) Interest is a fee charged by banks independent of market forces; D) Interest rates fluctuate based on the phases of the moon; E) Interest is a tax on savings that discourages postponing consumption
\\ \textit{Note:} The correct answer reflects the neoclassical view that individuals prefer to consume goods and services sooner rather than later, leading to the notion that interest serves as a compensation for delaying consumption, effectively acting as a penalty for those who choose to consume immediately.
\item \textbf{Answer:} A
\\ \textit{Options:} A) \$599.69; B) \$744.40; C) \$681.47; D) \$880
\\ \textit{Note:} The correct answer of $599.69 represents the present value of $1,000 adjusted for 12\% annual inflation over 4 years, calculated using the formula for real value, which accounts for the decrease in purchasing power due to inflation.
\end{enumerate}

\section*{True / False}
\begin{enumerate}[label=\arabic*.]
\setcounter{enumi}{5}
\item \textbf{Answer:} False
\\ \textit{Options:} A) True; B) False
\\ \textit{Note:} The economies and diseconomies of scale are primarily analyzed in the long-run context of production, as they involve changes in production capacity and efficiency that occur as a firm adjusts all its inputs.
\item \textbf{Answer:} False
\\ \textit{Options:} A) True; B) False
\\ \textit{Note:} A price floor is typically set above the equilibrium price to prevent prices from falling too low, which can lead to excess supply and reduced consumer spending.
\item \textbf{Answer:} True
\\ \textit{Options:} A) True; B) False
\\ \textit{Note:} GDP is indeed calculated annually by the Bureau of Economic Analysis to provide a comprehensive assessment of a country's economic performance by measuring the total value of all goods and services produced within its borders.
\item \textbf{Answer:} True
\\ \textit{Options:} A) True; B) False
\\ \textit{Note:} In a standard unrestricted tri-variate VAR(4) model, each of the three variables has four lags, resulting in 3 variables × 4 lags = 12 parameters, and since each variable has two other variables influencing it, the total number of parameters for the equations is 12 - 6 (one for each variable's own lagged value) = 6 parameters, excluding intercepts.
\item \textbf{Answer:} False
\\ \textit{Options:} A) True; B) False
\\ \textit{Note:} A cross-price elasticity of +2.0 indicates that goods X and Y are substitutes, as their demand moves in the same direction when the price of one changes, rather than being complementary goods which would have a negative cross-price elasticity.
\end{enumerate}

\section*{Long Form Response}
\begin{enumerate}[label=\arabic*.]
\setcounter{enumi}{10}
\item \textbf{Answer:} When both supply decreases and demand increases in the market for roses, the immediate effect is an upward pressure on price. The decrease in supply indicates a reduction in the availability of roses, while the increase in demand reflects a heightened desire for them. Consequently, the price of roses is likely to rise due to the scarcity created by the supply reduction coupled with the increased consumer interest. However, the impact on quantity is ambiguous; while the reduced supply would typically lower quantity sold, the increased demand could counteract this effect by encouraging higher sales. Thus, while we can confidently assert that prices will rise, the overall change in quantity remains uncertain as it depends on the relative magnitudes of the shifts in supply and demand.
\\ \textit{Note:} correct because a simultaneous decrease in supply and increase in demand leads to a higher equilibrium price due to scarcity, while the effect on equilibrium quantity remains uncertain as it depends on the relative magnitudes of the shifts in supply and demand.
\item \textbf{Answer:} In a perfectly competitive market, firms make labor hiring decisions based on the marginal cost (MC) of hiring additional labor relative to the marginal revenue product (MRP) of that labor. The MC of hiring labor, referred to as marginal resource cost (MRC), must equal the MRP for a firm to maximize profits. This implies that a firm will continue to hire more workers until the wage paid, which corresponds to the MRC, equals the MRP of the last worker hired. Thus, in equilibrium, when MRC = MRP, the wage reflects the additional value generated by the last unit of labor, ensuring efficient resource allocation and optimal labor employment. This relationship underscores the principle that firms in competitive markets rationally adjust their labor input to align with productivity and cost considerations.
\\ \textit{Note:} correct because it accurately describes how firms in a perfectly competitive market optimize labor hiring by equating the marginal resource cost (MRC) of hiring additional labor with the marginal revenue product (MRP) generated by that labor, ensuring profit maximization at the point where MRC equals MRP.
\end{enumerate}

\vfill
\noindent\textit{End of Answer Key}
\end{document}